\section{Introduzione}
\label{sec:introduzione_st}
Il seguente documento definisce la specifica tecnica del progetto Artificial QI, un software progettato per valutare l'affidabilità 
delle risposte fornite da un \glossario{LLM} rispetto a un dataset di domande.  L'obiettivo di questa specifica è fornire una 
definizione chiara e dettagliata delle scelte fatte dal team di sviluppo riguardo l'architettura del sistema, i design pattern 
utilizzati, le tecnologie adottate e le scelte progettuali effettuate.
Questo documento andrà a indicare:
\begin{itemize}
    \item Le architetture e i design pattern scelti sia lato front-end che back-end.
    \item Le tecnologie utilizzate per la realizzazione del software.
    \item La progettazione ad alto livello dei componenti sia lato front-end che back-end.
    \item La progettazione di dettaglio dei componenti sia lato front-end che back-end.
\end{itemize}



\subsection{Riferimenti}
\subsubsection{Normativi}
\begin{itemize}
    \item Norme di progetto
    \item Capitolato \url{https://www.math.unipd.it/~tullio/IS-1/2024/Progetto/C1.pdf}
    \item Regolamento del progetto: \url{https://www.math.unipd.it/~tullio/IS-1/2024/Dispense/PD1.pdf}
\end{itemize}
\subsubsection{Informativi}
\begin{itemize}
    \item Analisi dei Requisiti
    \item Glossario
    \item Pattern architetturali \url{https://www.math.unipd.it/~rcardin/swea/2022/Software%20Architecture%20Patterns.pdf}
    \item Architettura esagonale \url{https://github.com/rcardin/hexagonal}
    \item Progettazione software \url{https://www.math.unipd.it/~tullio/IS-1/2024/Dispense/T06.pdf}
    \item Analisi statica \url{https://www.math.unipd.it/~tullio/IS-1/2024/Dispense/T10.pdf}
    \item Analisi dinamica \url{https://www.math.unipd.it/~tullio/IS-1/2024/Dispense/T11.pdf}
    \item Pattern creazionali \url{https://www.math.unipd.it/~rcardin/swea/2022/Design%20Pattern%20Creazionali.pdf}
    \item Pattern strutturali \url{https://www.math.unipd.it/~rcardin/swea/2022/Design%20Pattern%20Strutturali.pdf}
    \item Pattern comportamentali \url{https://drive.google.com/file/d/1cpi6rORMxFtC91nI6_sPrG1Xn-28z8eI/view}
    \item Verbali interni
        \begin{itemize}
            \item 2024\_03\_24
        \end{itemize}
    \item Verbali esterni
\end{itemize}
